\section{Introduction}\label{sec:intro}

Chain-of-thought (CoT) prompting has emerged as the dominant technique for
eliciting complex reasoning from large language models (LLMs)
\cite{wei2022chain}. By decomposing a problem into sequential reasoning steps,
CoT enables LLMs to solve tasks that are intractable in a single forward pass.
Despite the empirical success of this paradigm, a rigorous mathematical
foundation for understanding \emph{when} and \emph{why} iterative reasoning
converges to correct solutions remains undeveloped.

Recent work has approached this question from several directions. Kim, Wu, Lee,
and Suzuki \cite{kim2025metastable} model CoT as a Markov chain on a finite
state space of logical assertions with metastable cluster structure,
establishing power-law hitting time bounds of order $\widetilde{\Theta}(KM/\varepsilon)$
for reaching target reasoning states. Carson \cite{carson2025stochastic}
develops a stochastic differential equation framework in which the reasoning
process is governed by a Fokker--Planck equation, identifying phase transitions
between convergent and divergent regimes. Bai, Kolter, and Koltun
\cite{bai2019deep} introduce deep equilibrium models, which formalize neural
network inference as convergence to a fixed point, with convergence rate
controlled by the spectral radius of the Jacobian. Ke et al.\
\cite{ke2024fixed} extend this to looped neural networks, proving exponential
convergence $\|x^{(L)} - p\| \le K^L \cdot c$ with contraction constant $K < 1$.

These results suggest two qualitatively different convergence behaviors:
exponential convergence under contraction conditions and slow convergence in the
presence of metastable traps. However, several fundamental questions remain open.
First, no prior work provides convergence analysis for CoT iterations in the
Wasserstein metric on probability distributions, despite Wasserstein distances
being the natural metric for distributional convergence of stochastic processes.
Second, the spectral theory of the CoT transition operator---connecting
eigenvalues to convergence rates---has not been developed in a unified manner.
Third, the conditions separating exponential from slow convergence regimes have
not been characterized in terms of problem parameters.

In this paper, we address these gaps by developing a rigorous mathematical
framework for CoT reasoning as a discrete dynamical system on probability
distributions. Our main contributions are as follows.

\begin{itemize}[leftmargin=*,itemsep=2pt,topsep=4pt]
\item We formalize CoT reasoning as a Markov transition operator
  $T \colon \mathcal{P}(S) \to \mathcal{P}(S)$ on the probability simplex
  equipped with the Wasserstein metric, and prove well-posedness and continuity
  (Lemma~\ref{lem:wellposed}).

\item We establish exponential convergence in $W_1$ under Dobrushin contraction,
  with rate governed by a Wasserstein contraction coefficient that is finer than
  the classical Dobrushin coefficient (Theorem~\ref{thm:exponential}).

\item We derive an exact first-order formula for the spectral gap of
  $(K,\varepsilon)$-metastable transition matrices, $\gamma = K\varepsilon/(K-1)$,
  via perturbation analysis of the aggregated inter-cluster chain
  (Theorem~\ref{thm:metastable}).

\item We characterize a phase transition at critical coupling
  $\varepsilon_c = (K-1)\gamma_{\mathrm{intra}}/K$ separating slow-mixing and
  fast-mixing regimes (Theorem~\ref{thm:phase}).

\item We verify all theoretical predictions computationally, achieving
  $r = 0.979$ correlation between predicted and empirical convergence rates.
\end{itemize}

The paper is organized as follows. Section~\ref{sec:prelim} introduces the
state space, Wasserstein metrics, transition operators, and contraction
coefficients. Section~\ref{sec:results} contains the main results with
complete proofs. Section~\ref{sec:discussion} discusses implications for CoT
prompting, connections to prior work, and open questions.
Section~\ref{sec:conclusion} summarizes our contributions.
